% Part: first-order-logic
% Chapter: introduction
% Section: sentences

\documentclass[../../../include/open-logic-section]{subfiles}

\begin{document}

\olfileid{fol}{int}{snt}

\olsection{\usetoken{P}{sentence}}

بسیار خوب، اکنون طرحی از تعریف ارضا («صادق در») برای !!{structure}s و
!!{formula}s داریم. اما این تعریف به جزء دیگری---!!a{variable}
گمارش---نیاز دارد، حال آنکه خواستهٔ ما تعریف !!{sentence}s بود. چگونه
از گمارش‌ها رها شویم و !!{sentence}s چه هستند؟

احتمالاً از نخستین آشنایی خود با منطق صوری، بحث رابطهٔ میان
!!{variable}s و سورها را به یاد دارید. پس از هر سور همیشه
!!a{variable} می‌آید و در بخشی از !!{sentence} که سور بر آن اعمال
می‌شود («دامنهٔ» سور)، می‌گوییم !!{variable} به‌وسیلهٔ آن سور «مقید»
شده است. در !!{formula}s لازم نبود هر !!{variable} سوری متناظر داشته
باشد؛ !!{variable}s بدون سور متناظر «آزاد» یا «نامقید»اند. همهٔ
!!{formula}s بدون !!{variable}s آزاد را !!{sentence} خواهیم گرفت.

باز هم فهم ایدهٔ شهودیِ اینکه یک رخداد !!a{variable} در
!!a{formula}~$!A$ چه هنگام مقید است، کدام سور آن را مقید می‌کند و چه
هنگام آزاد است، دشوار نیست. شاید روشی برای آزمودن این امر آموخته باشید
که مثلاً شامل شمردن پرانتزها باشد. اما باید بر تعریفی دقیق اصرار کنیم؛
و چون !!{formula}s را به‌طور استقرایی تعریف کرده‌ایم، می‌توانیم رخدادهای
آزاد و مقید !!a{variable}~$x$ در !!a{formula}~$!A$ را نیز با استقرا
تعریف کنیم. برای نمونه، در زبان ساده‌شدهٔ ما تعریف می‌تواند چنین باشد:
\begin{enumerate}
  \item اگر $!A$ اتمی باشد، همهٔ رخدادهای $x$ در آن آزادند (یعنی رخداد
  $x$ در~$\Atom{\Obj P}{x}$ آزاد است).
  \item اگر $!A$ به صورت $\lnot !B$ باشد، یک رخداد~$x$ در~$\lnot !B$
  آزاد است هرگاه و تنها هرگاه رخداد متناظر~$x$ در~$!B$ آزاد باشد (یعنی
  رخدادهای آزاد متغیرها در~$!B$ دقیقاً رخدادهای متناظر در~$\lnot !B$
  هستند).
  \item اگر $!A$ به صورت $(!B \land !C)$ باشد، یک رخداد~$x$ در
  ~$(!B \land !C)$ آزاد است هرگاه و تنها هرگاه رخداد متناظر~$x$ در
  ~$!B$ یا~$!C$ آزاد باشد.
  \item اگر $!A$ به صورت $\lexists[x][!B]$ باشد، هیچ رخداد $x$ در~$!A$
  آزاد نیست؛ اگر به صورت $\lexists[y][!B]$ باشد که $y$ !!{variable}ی
  متفاوت از~$x$ است، یک رخداد~$x$ در $\lexists[y][!B]$ آزاد است هرگاه
  و تنها هرگاه رخداد متناظر~$x$ در~$!B$ آزاد باشد.
\end{enumerate}

هنگامی که تعریفی دقیق از رخدادهای آزاد و مقید متغیرها داشته باشیم،
به‌سادگی می‌توانیم بگوییم: هر !!{formula} بدون رخداد آزاد
!!{variable}s، !!a{sentence} است.

\end{document}
